% ---------------------------------------
% -------------- Preamble ---------------
% ---------------------------------------
\documentclass[a4paper,12pt,twoside,numbers=noenddot]{scrartcl}

% Dokument-bezogene Packages
\usepackage[utf8]{inputenc}
%\usepackage[T1]{fontenc}
\usepackage[ngerman]{babel}
\usepackage{setspace}
\onehalfspacing
\setstretch{1.2}								% Zeilenabstandsfaktor

% Kopf- und Fußzeilen, Ränder
\setlength{\headheight}{2\baselineskip}			% Kopfzeile breiter für 2 Zeilen
\usepackage[a4paper,bottom=3.5cm,top=3cm]{geometry}
\usepackage{scrpage2}

% Abschnitts-Überschriften/-Styles
\renewcommand*\sectfont{\normalfont\rmfamily\bfseries}

% Text-Darstellung
\usepackage[protrusion=true,expansion=true]{microtype}

% Mathe-Pakete
\usepackage{amsmath}
\usepackage{amsfonts}
\usepackage{amssymb}

% Links
\usepackage{hyperref}

% Enumerationen
\usepackage{enumerate}

% Theoreme
\usepackage[amsmath,hyperref,thmmarks]{ntheorem}
\theoremseparator{.}
\theorembodyfont{\normalfont}
\newtheorem{definition}{Definition}[subsection]	% Definition
\newtheorem{lemma}{Lemma}[subsection]			% Lemma
\newtheorem{theorem}{Theorem}[subsection]		% Theorem
\newtheorem{example}{Beispiel}					% Beispiel
\theoremstyle{nonumberplain}
\newtheorem{remark}{Bemerkung}					% Bemerkung
\theoremsymbol{\ensuremath{\square}}
\newtheorem{proof}{Beweis}						% Beweis

% Grafiken
\usepackage{graphicx}
\graphicspath{{./images/}}
\usepackage{tikz}
\usetikzlibrary{arrows}
\usetikzlibrary{patterns}
\usepackage{pgfplots} % Plots
\pgfplotsset{compat=1.7}

% Tabellen
\usepackage{booktabs}
\usepackage{tabu}

% Algorithmen
\usepackage[ruled,noline,commentsnumbered,linesnumbered,noend,german]{algorithm2e}
\let\chapter\section % Kompatibilität mit KOMA

% Literatur
\usepackage[square]{natbib}
\renewcommand*{\bibfont}{\interlinepenalty 10000\relax}	% Kein Seitenumbruch innerhalb von Bib-Eintrag

% Eigene Befehle
\newcommand{\absatz}{\\[\baselineskip]}
\newcommand{\emptypage}{\clearpage\thispagestyle{empty}\mbox{}\newpage}
\newcommand{\emptypagenumbered}{\clearpage\thispagestyle{scrheadings}\mbox{}\newpage}

% Sonstige Konfiguration
\usepackage{etoolbox}
\let\stdsection\section
\renewcommand\section{\newpage\stdsection}		% Neue Seite pro Kapitel

% Flag für Erklärung
\newtoggle{erklaerung}
\toggletrue{erklaerung}
%\togglefalse{erklaerung}

% ---------------------------------------
% -------------- Document ---------------
% ---------------------------------------
\begin{document}

\lefthyphenmin=3								% Nicht weniger als 3 Buchstaben trennen

\pagenumbering{roman}
% Titel-Seite
\begin{titlepage}
\newgeometry{left=3.5cm,right=2.5cm,bottom=3.5cm,top=3cm}
\pagestyle{empty}
\definecolor{TUGreen}{rgb}{0.517,0.721,0.094}
\vspace*{-2cm}
\newlength{\links}
\setlength{\links}{-1.5cm}
%\sffamily
\rmfamily
\hspace*{\links}
\begin{minipage}{12.5cm}
\includegraphics[width=8cm]{tud_logo_rgb}
%\hspace*{-0.25cm} \textbf{TECHNISCHE UNIVERSIT"AT DORTMUND}\\
%\hspace*{-1.2cm} \rule{5mm}{5mm} \hspace*{0.1cm} FACHBEREICH INFORMATIK\\
\end{minipage}

\vspace*{4cm}

\hspace*{\links}
\hspace*{-0.2cm}
\begin{minipage}{9cm}
\large
\begin{center}
{\Large Masterarbeit} \\
\vspace*{1cm}
\textbf{Kernel k-means Methoden zur spektralen Clusteranalyse von Graphen} \\
\vspace*{1cm}
Lukas Pradel\\
% \vspace*{1cm}
\today
\end{center}
\end{minipage}
\normalsize
%\vspace*{5.5cm}
\vspace*{4.8cm}

% \hspace*{\links}

\vspace*{2.1cm}

\hspace*{\links}
\begin{minipage}[b]{5cm}
% \normalsize
\raggedright
Gutachter: \\
Prof. Dr. Christian Sohler \\
Dipl.-Inf. Melanie Schmidt \\
\end{minipage}

\vspace*{2.5cm}
\hspace*{\links}
%\begin{minipage}[b]{8cm}
\begin{minipage}[b]{11cm}
% \normalsize
\raggedright
Technische Universit"at Dortmund \\
Fakult"at f"ur Informatik\\
Lehrstuhl II - Effiziente Algorithmen und Komplexitätstheorie\\
\href{http://ls2-www.cs.tu-dortmund.de/}{http://ls2-www.cs.tu-dortmund.de/}
\end{minipage}
\end{titlepage}
\restoregeometry

% Danksagung
\thispagestyle{empty}
\section*{}
\vfill
Hier kommt eine Danksagung hin!
\emptypagenumbered
\emptypagenumbered

% Inhaltsverzeichnis
%\pagestyle{empty}
\pagestyle{scrheadings}
\tableofcontents
\emptypagenumbered

% Page-Style für Kapitel
\pagestyle{scrheadings}

%\automark[subsection]{section}

\pagenumbering{arabic}
%\setcounter{page}{1}	% hier neu zählen

% Kapitel

% Anhänge

% Literatur
\emptypagenumbered
\addcontentsline{toc}{section}{Literatur}
%\bibliographystyle{ba}
\bibliography{bib}
\nocite{*}

% Erklärung
\iftoggle{erklaerung}{
	\emptypagenumbered
	\emptypagenumbered
	%\markboth{}{ERKLÄRUNG}
	\addcontentsline{toc}{section}{Erklärung}
	Hiermit versichere ich, dass ich die vorliegende Arbeit selbstständig verfasst habe und keine
anderen als die angegebenen Quellen und Hilfsmittel verwendet sowie Zitate kenntlich gemacht habe.\\\\
Dortmund, den \today \\\\\\\\
Lukas Pradel
}

\emptypagenumbered

\end{document}